\section{问题重述}\label{sec:problem_restatement}

多模态情感识别利用文本、语音和视觉观测判断多模态视频片段的情感状态\cite{baltrusaitis2019multimodal,poria2017affective}。三类观测的来源与时间粒度不同，最终又需要同时给出类别、连续强度及预测依据。本题依次提供原始视频、已对齐特征和待预测样本，要求将特征构建、缺失条件下的预测与预测解释放在同一研究任务中处理。

\textbf{问题一：原始特征的构建与时间统一。}附件1包含100条原始多模态视频。需要从每条视频分别获得文本、语音和视觉特征，同时记录这些特征对应的时间位置。由于词语、声学片段与画面不能按原生序号直接对应，还需将三路信息放到可比较的公共时间位置，形成统一长度的特征及有效性记录，并使结果能够回查到原始观测。

\textbf{问题二：完整与局部缺失输入下的情感预测。}附件2提供统一到50个时间位置的三模态特征及标签。模型需要同时预测离散的情感极性和连续的情感强度，并比较完整输入与局部连续缺失输入下的表现。缺失可涉及不同模态、不同时间位置和不同持续长度，因此评价应覆盖这些变化，而不能只考察一种固定遮挡。训练和评价完成后，还需对附件3的无标签样本给出两项预测。

\textbf{问题三：预测依据与局部位置的解释。}在问题二的预测模型上，需要判断各模态对当前类别和强度输出分别起到什么作用，考察两种模态同时输入时的作用关系，并定位对输出影响较大的连续位置。解释结果应能对应到模型实际使用的输入；对附件4的无标签样本，还需同时报告预测和相应解释。

三个问题分别围绕多模态特征构建、连续缺失条件下的情感预测和预测解释展开。其中，问题一与问题二分别基于附件1和附件2开展，数据来源相互独立；问题三则沿用问题二建立的预测模型，进一步分析各模态及局部输入对预测结果的作用。
